% Ch. 29 : le cycle de vie d'un modèle, deux boucles
\documentclass[border=6pt]{standalone}
\usepackage{coursfig}
\begin{document}
\begin{tikzpicture}[x=1mm,y=1mm]
  \tikzset{st/.style={box=#1, text width=30mm, minimum height=13mm}}
  % boucle d'expérimentation (gauche)
  \node[st=Ocre] (d) at (0,24)  {\textbf{données}\sub{collecte, instantané}};
  \node[st=Ocre] (v) at (0,0)   {\textbf{validation}\sub{des données}};
  \node[st=Ocre] (f) at (0,-24) {\textbf{préparation}\sub{variables}};
  \node[st=Enc]  (t) at (44,-24){\textbf{entraînement}};
  \node[st=Enc]  (e) at (44,0)  {\textbf{évaluation}\sub{contre le champion}};
  \draw[flow] (d) -- (v); \draw[flow] (v) -- (f); \draw[flow] (f) -- (t); \draw[flow] (t) -- (e);
  \draw[dflow=Muted, rounded corners=4pt] (e.north) |- node[elabel, pos=.75, above=1mm, fill=none]{nouvelle idée} (d.east);
  % boucle de production (droite)
  \node[st=Sea]  (p) at (114,0)  {\textbf{promotion}\sub{registre de modèles}};
  \node[st=Sea]  (s) at (158,0)  {\textbf{service}\sub{prédictions}};
  \node[st=Plum] (m) at (158,-24){\textbf{surveillance}\sub{qualité, dérive}};
  \node[box=Brick, text width=30mm, minimum height=13mm] (r) at (114,-24) {\textbf{déclencheur}\sub{dérive, calendrier}};
  \draw[flow=Sea, line width=1pt] (e) -- node[elabel, above=1mm, fill=none]{meilleur ?} (p);
  \draw[flow] (p) -- (s); \draw[flow] (s) -- (m); \draw[flow] (m) -- (r);
  \draw[flow=Brick, rounded corners=4pt] (r.south) -- ++(0,-8) -| node[elabel, pos=.25, below=1mm, fill=none, text=Brick]{réentraînement} (f.south);
  \begin{scope}[on background layer]
    \node[dframe=Ocre, fit=(d)(t)(e), inner sep=5mm, yshift=3mm, inner ysep=8mm] (b1) {};
    \node[dframe=Sea, fit=(p)(s)(m)(r), inner sep=5mm, yshift=3mm, inner ysep=8mm] (b2) {};
  \end{scope}
  \node[ftitle, text=Ocre, anchor=north west] at (b1.north west) {EXPÉRIMENTATION};
  \node[ftitle, text=Sea, anchor=north east] at (b2.north east) {PRODUCTION};
\end{tikzpicture}
\end{document}
